\documentclass[11pt]{article}
\usepackage{amsmath,amssymb,amsthm,fullpage}
\theoremstyle{plain}
\newtheorem{theorem}{Theorem}
\newtheorem{proposition}[theorem]{Proposition}
\newtheorem{lemma}[theorem]{Lemma}
\newtheorem{corollary}[theorem]{Corollary}
\theoremstyle{definition}
\newtheorem{definition}[theorem]{Definition}
\newtheorem{remark}[theorem]{Remark}
\newtheorem{gap}[theorem]{Gap}
\newcommand{\SSYT}{\mathrm{SSYT}}
\newcommand{\Z}{\mathbb{Z}}
\newcommand{\charge}{\mathrm{c}}
\DeclareMathOperator{\dom}{\trianglerighteq}

\title{Charge is never constant on a nontrivial fibre:\\
       the monomial locus of the Kostka--Foulkes polynomials}
\author{Clio}
\date{9 October 2026}

\begin{document}
\maketitle

\begin{abstract}
Question Q380 asked for the pairs $(\lambda,\mu)$ with $\#\SSYT(\lambda,\mu)\ge 2$ for
which the Kostka--Foulkes polynomial $K_{\lambda\mu}(t)$ is a monomial --- equivalently,
the pairs on which the Lascoux--Schützenberger charge statistic is constant on a fibre
with something to be constant \emph{on}.  \textbf{The answer is that there are none},
and the reason is one line: $K_{\lambda\mu}(t)$ is \emph{monic}, so a monomial is forced
to be $t^{n(\mu)-n(\lambda)}$ with coefficient $1$, and its value at $t=1$ --- the Kostka
number --- is therefore $1$.  Hence
\[
   K_{\lambda\mu}(t)\ \text{is a monomial}
   \iff \#\SSYT(\lambda,\mu)=1
   \iff K_{\lambda\mu}(t)=t^{\,n(\mu)-n(\lambda)} .
\]
The monomial locus coincides exactly with the Kostka-number-one locus, so Q380's
``boundary'' is the boundary of $\{K_{\lambda\mu}=1\}$, and \S\ref{sec:boundary} measures
it.  Three things are proved there: (i) for fixed $\lambda$ the locus is an \emph{order
filter} in the dominance order, hence determined by its minimal elements, of which there
are \textbf{at most two} for every $\lambda$ with $|\lambda|\le 10$; (ii) two exact
reductions $R_1,R_2$ account for all but $10$ of the $1044$ irreducible dominance pairs
with $|\lambda|\le 10$; (iii) no criterion depending only on the difference vector
$\lambda-\mu$, and none depending only on the dominance-slack vector, can decide the
question --- both are refuted by explicit witness pairs.  A closed-form characterisation
of $\{K_{\lambda\mu}=1\}$ is \textbf{not} obtained and the obstruction is stated
precisely in \S\ref{sec:gaps}.

Everything is verified by two mechanism-disjoint engines and against Macdonald's printed
tables, on an exhaustive census of all $1719$ nonempty dominance pairs with
$|\lambda|\le 10$, of which $1245$ are non-vacuous ($\#\SSYT\ge 2$).  The consequence for
the migration no-go of 5 October, which is what Q380 was gating, is
Corollary~\ref{cor:nogo}: \emph{the no-go's conclusion has no non-vacuous instance for
charge at all.}
\end{abstract}

\paragraph{Attribution, stated first.}  Theorem~\ref{thm:main} is a two-line deduction
from Macdonald III (6.5), and I expect it is folklore; I have not searched for a prior
statement and make no priority claim.  Its value here is use-value: it is the fact that
Q370 was waiting on, and the fact that stops a future session applying the 5 October
no-go where it has no content.  The proof of monicity reproduced in
Lemma~\ref{lem:monic} is Macdonald's, from III \S6 Example 4; it is reproduced rather
than cited because the whole paper rests on it.  What is new here, as far as I can tell,
is only the boundary measurement of \S\ref{sec:boundary} and the two obstruction
statements of \S\ref{sec:obstruction}.

\section{Notation and the two inputs}

Partitions $\lambda,\mu$ of $N$; $\lambda' $ is the conjugate; $\ell(\lambda)$ the
number of parts; $n(\lambda)=\sum_{i\ge1}(i-1)\lambda_i$ (Macdonald I (1.5)).  Dominance:
$\lambda\dom\mu$ iff $\sum_{i\le k}\lambda_i\ge\sum_{i\le k}\mu_i$ for all $k$.  We write
$\SSYT(\lambda,\mu)$ for the set of semistandard Young tableaux of shape $\lambda$ and
content $\mu$, and $K_{\lambda\mu}=\#\SSYT(\lambda,\mu)$ for the Kostka number.  The
Kostka--Foulkes polynomials $K_{\lambda\mu}(t)$ are the entries of the transition matrix
$K(t)=M(s,P)$ from Schur to Hall--Littlewood functions,
$s_\lambda=\sum_\mu K_{\lambda\mu}(t)P_\mu(x;t)$.

Two facts are taken from Macdonald, \emph{Symmetric Functions and Hall Polynomials}, 2nd
ed., III \S6 (pp.\ 238--245), deep-read 9 October 2026, with pp.\ 239--240 read off page
images rendered at 200\,dpi rather than off the text layer.

\begin{lemma}[Macdonald III (6.2)]\label{lem:at1}
$K(1)$ is the Kostka matrix: $K_{\lambda\mu}(1)=K_{\lambda\mu}=\#\SSYT(\lambda,\mu)$.
Also $K_{\lambda\mu}(t)=0$ unless $\lambda\dom\mu$, and $K_{\lambda\lambda}(t)=1$.
\end{lemma}

\begin{lemma}[Macdonald III (6.5)(ii)]\label{lem:monic}
If $\lambda\dom\mu$ then $K_{\lambda\mu}(t)$ is \emph{monic} of degree
$n(\mu)-n(\lambda)$.
\end{lemma}

\begin{proof}[Proof (Macdonald III \S6, Example 4), reproduced]
Let $n\ge\ell(\mu)$, let $\varepsilon_1,\dots,\varepsilon_n$ be the standard basis of
$\Z^n$, $R^+=\{\varepsilon_i-\varepsilon_j: i<j\}$, and
$\delta=(n-1,n-2,\dots,1,0)$.  For $\xi\in\Z^n$ with $\sum\xi_i=0$ set
\[
   P(\xi;t)=\sum_{(m_\alpha)} t^{\sum_\alpha m_\alpha},
\]
summed over families $(m_\alpha)_{\alpha\in R^+}$ of non-negative integers with
$\xi=\sum_\alpha m_\alpha\alpha$; this is the Lusztig $t$-analogue of Kostant's partition
function.  Writing $\xi=\sum_{i<n}\eta_i(\varepsilon_i-\varepsilon_{i+1})$ with
$\eta_i=\xi_1+\dots+\xi_i$, one has $P(\xi;t)\ne0$ iff all $\eta_i\ge 0$, and then
$P(\xi;t)$ is monic of degree $\sum_i\eta_i=\langle\xi,\delta\rangle$.  (Monic: the unique
family of maximal total weight is $m_{\varepsilon_i-\varepsilon_{i+1}}=\eta_i$ and all
other $m_\alpha=0$, because replacing a simple-root unit by a longer root strictly lowers
$\sum m_\alpha$.)

By III (6.3), $K(t)'$ is the matrix of $\prod_{i<j}(1-tR_{ij})^{-1}$, so $K_{\lambda\mu}(t)$
is the coefficient of $x^{\lambda+\delta}$ in
$\sum_{w\in S_n}\varepsilon(w)\,w\!\left(x^{\mu+\delta}\prod_{i<j}(1-t x_i^{-1}x_j)^{-1}\right)$,
whence
\begin{equation}\label{eq:alt}
   K_{\lambda\mu}(t)=\sum_{w\in S_n}\varepsilon(w)\,
      P\!\left(w^{-1}(\lambda+\delta)-(\mu+\delta);\,t\right).
\end{equation}
For the degree of the $w$-term,
\[
  \big\langle w^{-1}(\lambda+\delta)-(\mu+\delta),\ \delta\big\rangle
   =\langle \lambda+\delta,\,w\delta\rangle-\langle\mu+\delta,\delta\rangle
   \le \langle \lambda+\delta,\,\delta\rangle-\langle\mu+\delta,\delta\rangle
   = n(\mu)-n(\lambda),
\]
using $\sum_i(n-i)\lambda_i=(n-1)|\lambda|-n(\lambda)$.  The inequality is the
rearrangement inequality applied to the two strictly decreasing vectors $\lambda+\delta$
and $\delta$, and it is \emph{strict} for every $w\ne 1$.  The $w=1$ term is
$P(\lambda-\mu;t)$, which is nonzero precisely because $\lambda\dom\mu$, and is monic of
degree $n(\mu)-n(\lambda)$.  Every other term has strictly smaller degree, so no
cancellation can touch the leading coefficient.
\end{proof}

We will also use, but only in Corollary~\ref{cor:charge}, the Lascoux--Schützenberger
theorem, Macdonald III (6.5)(i): $K_{\lambda\mu}(t)=\sum_{T\in\SSYT(\lambda,\mu)}t^{\charge(T)}$,
$\charge$ the charge of the reading word.  \textbf{The main theorem does not depend on it.}

\section{The theorem}

\begin{theorem}\label{thm:main}
Let $\lambda\dom\mu$ be partitions of $N$.  The following are equivalent.
\begin{enumerate}
\item[(i)] $K_{\lambda\mu}(t)$ has exactly one nonzero coefficient.
\item[(ii)] $K_{\lambda\mu}(t)=t^{\,n(\mu)-n(\lambda)}$.
\item[(iii)] $\#\SSYT(\lambda,\mu)=1$.
\end{enumerate}
\end{theorem}

\begin{proof}
(i)$\Rightarrow$(ii).  Write $K_{\lambda\mu}(t)=c\,t^{d}$ with $c\neq 0$.  By
Lemma~\ref{lem:monic} the leading coefficient is $1$ and the degree is
$n(\mu)-n(\lambda)$; a single-term polynomial has its only coefficient as its leading
coefficient, so $c=1$ and $d=n(\mu)-n(\lambda)$.

(ii)$\Rightarrow$(iii).  Evaluate at $t=1$ and apply Lemma~\ref{lem:at1}:
$\#\SSYT(\lambda,\mu)=K_{\lambda\mu}(1)=1$.

(iii)$\Rightarrow$(i).  $K_{\lambda\mu}(t)$ has non-negative integer coefficients
(Lemma~\ref{lem:at1} with III (6.5)(i), or Lusztig's geometric interpretation, III \S6
Example 6) summing to $K_{\lambda\mu}(1)=1$, so exactly one coefficient is nonzero.
\end{proof}

\begin{corollary}[the answer to Q380]\label{cor:q380}
$\big\{(\lambda,\mu):\ \#\SSYT(\lambda,\mu)\ge 2\ \text{and}\ K_{\lambda\mu}(t)
\ \text{is a monomial}\big\}=\varnothing.$
\end{corollary}

\begin{corollary}[charge separates]\label{cor:charge}
If $\#\SSYT(\lambda,\mu)\ge 2$ then the charge statistic takes at least two distinct
values on $\SSYT(\lambda,\mu)$.  More precisely, $\max_T \charge(T)=n(\mu)-n(\lambda)$ and
it is attained by exactly one tableau.
\end{corollary}

\begin{proof}
By III (6.5)(i) the multiset of charges is the coefficient list of $K_{\lambda\mu}(t)$;
by Lemma~\ref{lem:monic} the top coefficient is $1$ and the top degree is
$n(\mu)-n(\lambda)$, so the maximum charge is attained once.  If
$\#\SSYT(\lambda,\mu)\ge2$, some other tableau has strictly smaller charge.
\end{proof}

\begin{remark}
Corollary~\ref{cor:charge} is the statement with content.  Theorem~\ref{thm:main} says
the \emph{only} way for charge to be constant on a fibre is for the fibre to have nothing
to be constant on --- which is exactly the degenerate situation that Q380 was constructed
to exclude.  Equivalently: \emph{the unique maximal-charge tableau is the entire
obstruction}, and it exists in every fibre.
\end{remark}

\subsection{What this does to the migration no-go, and to Q370}

\begin{corollary}[the no-go is vacuous for charge]\label{cor:nogo}
Let $\phi$ be any statistic on $\SSYT(\lambda,\mu)$ and suppose the no-go of
\texttt{2026-10-05-c3-migration-length-grading.tex} applies to it --- i.e.\ $\phi$ is a
potential difference for some monotone process on the fibre, so that the fibre generating
function $\sum_T q^{\phi(T)}$ is a monomial $c\,q^{\ell_0}$.  Taking $\phi=\charge$, that
conclusion holds \emph{if and only if} $c=1$.  Therefore the no-go, applied to the charge
statistic on Kostka--Foulkes fibres, has \textbf{no non-vacuous instance whatsoever}:
its hypothesis forces precisely the fibres on which its conclusion is empty.
\end{corollary}

\begin{proof}
Immediate from Theorem~\ref{thm:main}: ``$\sum_T q^{\charge(T)}$ is a monomial'' is
equivalent to $\#\SSYT(\lambda,\mu)=1$, which is $c=1$.
\end{proof}

This settles the gating question behind Q370 (\emph{why does charge evade the potential
no-go: by reversibility, by branching, or because $K_{\lambda\mu}(t)$ is not a fibre
generating function of a monotone process at all?}) in the following sense, and only in
this sense: \textbf{charge does not need to evade the no-go, because the no-go has no
purchase on it.}  Whenever the fibre has $\ge 2$ elements, Corollary~\ref{cor:charge}
says $\sum_T q^{\charge(T)}$ is not a monomial, so by the contrapositive of the no-go,
\emph{charge is not a potential difference for any monotone process on that fibre}.  The
third horn of Q370 is the right one, and it is now a theorem rather than an assertion.
The amendment owed to the 5 October artifact (its \textbf{[A3]}) is therefore: on the
Kostka--Foulkes side the $c=1$ degeneracy is not an accident of the sample, it is forced.

\section{Engines, controls, and the census}\label{sec:census}

\paragraph{Engine A (charge).}  $\SSYT(\lambda,\mu)$ built by Gelfand--Tsetlin
horizontal-strip growth; each tableau then verified \emph{entrywise} for shape, content,
weak rows and strict columns; $K_{\lambda\mu}(t)=\sum_T t^{\charge(T)}$ with charge
computed by the Lascoux--Schützenberger standard-subword rule.

\paragraph{Engine B (Weyl alternating sum).}  Equation~\eqref{eq:alt}: a signed sum over
$S_n$ of values of the $t$-analogue of Kostant's partition function, the latter computed
by exhaustive enumeration of non-negative integer families $(m_\alpha)_{\alpha\in R^+}$.
No tableau, no reading word and no charge occurs anywhere in Engine~B.  The two engines
therefore differ in mechanism, not merely in basis.

\paragraph{Engine C (published).}  Macdonald's printed $K(t)$ matrices for $n\le 5$
(III \S6, p.\ 239), read off a 200\,dpi page image.

\begin{center}
\begin{tabular}{lll}
\hline
check & scope & result\\
\hline
A vs.\ B & all $53$ dominance pairs, $N\le5$ & $0$ mismatches\\
A vs.\ C & all printed entries, $n\le5$ & $0$ mismatches\\
\hline
\end{tabular}
\end{center}

\begin{remark}[the text layer lies]
The brief's warning was correct and it fired.  The OCR text layer of p.\ 239 column-shifts
the $n=5$ matrix, yielding \emph{$K_{(2,2,1),(2,1^3)}(t)=t$} --- a plausible-looking
monomial on a fibre that in fact has two tableaux and $K=t+t^2$.  Had that reading been
trusted, this paper would have opened with a counterexample to its own theorem.  It was
caught by hand before it was caught by the engines.  Pages that carry load are read as
images.
\end{remark}

\paragraph{Planted controls.}  Before reading any verdict:
\begin{itemize}
\item \emph{The statistic varies with the claim.}  The detector is
  $\#\{\text{nonzero coefficients}\}$.  On $(\lambda,\mu)=((3,1),(1^4))$ it reads $3$.
  Replacing $\charge$ by a constant function on that fibre makes it read $1$ with
  $\#\SSYT=3$, i.e.\ the planted pair \emph{does} enter the target set of
  Corollary~\ref{cor:q380}.  The detector is not a constant function of the thing it
  tests, and the empty target set below is a measurement, not a tautology.
\item \emph{The entrywise tableau check fires.}  Corrupting row order, a column
  inequality, or the content of a genuine tableau is rejected in all three cases.
\item \emph{Negative control.}  All $201$ pairs with $N\le 7$ failing dominance return
  $\SSYT=\varnothing$ (not a crash, not garbage).
\item \emph{Cardinality is not content.}  All $2767$ tableaux generated for $N\le8$ pass
  the entrywise check, and no generated set contains a duplicate.
\end{itemize}

\paragraph{The census.}  All pairs of partitions of $N$, $1\le N\le 10$.

\begin{center}
\begin{tabular}{rrrrrr}
\hline
$N$ & $\#$partitions & pairs & $\SSYT\neq\varnothing$ & $\#\SSYT\ge2$ & $K$ monomial\\
\hline
 1 &  1 &    1 &   1 &   0 &   1\\
 2 &  2 &    4 &   3 &   0 &   3\\
 3 &  3 &    9 &   6 &   1 &   5\\
 4 &  5 &   25 &  15 &   4 &  11\\
 5 &  7 &   49 &  28 &  12 &  16\\
 6 & 11 &  121 &  64 &  32 &  32\\
 7 & 15 &  225 & 116 &  72 &  44\\
 8 & 22 &  484 & 238 & 161 &  77\\
 9 & 30 &  900 & 430 & 320 & 110\\
10 & 42 & 1764 & 818 & 643 & 175\\
\hline
total & & 3582 & 1719 & 1245 & 474\\
\hline
\end{tabular}
\end{center}

The non-vacuous population is $1245$, so the range is not vacuous; it first becomes
nonempty at $N=3$.  Over all $1719$ nonempty pairs:
\[
 \#\{\deg K_{\lambda\mu}\ne n(\mu)-n(\lambda)\}=0,\qquad
 \#\{\text{leading coefficient}\ne1\}=0,
\]
\[
 \#\{\#\SSYT\ge2 \ \text{and}\ K_{\lambda\mu}(t)\ \text{a monomial}\}=\mathbf{0},
\]
and $474=1719-1245$ monomial pairs, i.e.\ the monomial count equals the $\#\SSYT=1$ count
at every $N$ separately.  This is Theorem~\ref{thm:main} and Lemma~\ref{lem:monic},
measured.

\paragraph{The two-row family (the brief's vacuity check).}  For all $N\le12$ and all
$j\le k$, $K_{(N-j,j),(N-k,k)}(t)=t^{\,k-j}$ with $\#\SSYT=1$: $138$ pairs, $0$ failures.
So \texttt{lem:mono} of \texttt{2026-10-07-two-part-green-polynomials.tex} is indeed a
\emph{singleton} fibre and is vacuous for Q380.  This is now also a corollary of
Proposition~\ref{prop:R2} below rather than a computation.

\section{The boundary}\label{sec:boundary}

By Theorem~\ref{thm:main} the monomial locus \emph{is} the Kostka-number-one locus
\[
  \mathcal{M}(\lambda)=\{\mu:\ K_{\lambda\mu}=1\}\subseteq\{\mu:\lambda\dom\mu\},
\]
so Q380's deliverable --- the boundary --- is the boundary of $\mathcal M(\lambda)$.

\subsection{$\mathcal M(\lambda)$ is an order filter}

\begin{proposition}\label{prop:filter}
If $\lambda\dom\nu\dom\mu$ then $K_{\lambda\nu}\le K_{\lambda\mu}$.  Consequently
$\mathcal M(\lambda)$ is an order filter (up-set) in the dominance interval
$[(1^N),\lambda]$, and is determined by its set of minimal elements.
\end{proposition}

\begin{proof}[Proof sketch, and the honest status]
This monotonicity of Kostka numbers in the dominance order on the content is classical.
Since $K_{\lambda\mu}$ depends only on the multiset of parts of $\mu$, it suffices to
treat $\nu=(\dots,a+1,b-1,\dots)$, $\mu=(\dots,a,b,\dots)$ with $a\ge b$ in two adjacent
positions $i,i+1$.  Restricting the crystal $B(\lambda)$ to the $i$-th $\mathfrak{sl}_2$
decomposes $\bigsqcup_c \SSYT(\lambda,\cdot)$ into $i$-strings, along each of which the
multiset of values of $m_i-m_{i+1}$ is symmetric about $0$ and consists of consecutive
integers of one parity; hence $\#\{m_i-m_{i+1}=d\}$ is weakly decreasing in $|d|$, which
is the claim.  \textbf{I have not written this argument out in full} --- the string
decomposition is standard but I am citing it, not proving it here.  Verified
computationally on all $6750$ triples $\lambda\dom\nu\dom\mu$ with $N\le9$: $0$
violations.
\end{proof}

\begin{theorem}[the measured boundary]\label{thm:minimal}
For every partition $\lambda$ with $|\lambda|\le 10$, the filter $\mathcal M(\lambda)$ has
\textbf{at most two} minimal elements.
\end{theorem}

This is a measurement over all $42+30+\dots=136$ partitions with $|\lambda|\le10$, not a
proof, and I do not conjecture that $2$ is a bound for all $N$ --- see \S\ref{sec:gaps}.
Examples at $N=10$ with exactly two minimal elements:
\[
\begin{array}{llll}
 (8,2)&\to\{(8,1,1),(5,5)\} &\quad (6,2,2)&\to\{(4,4,1,1),(4,3,3)\}\\
 (7,3)&\to\{(7,1^3),(5,5)\} &\quad (5,5)&\to\{(5,1^5),(4,3,3)\}\\
 (6,4)&\to\{(6,1^4),(5,5)\} &\quad (4,4,2)&\to\{(4,4,1,1),(4,3,3)\}\\
 (6,3,1)&\to\{(6,2,2),(5,4,1)\} &\quad (4,4,1,1)&\to\{(4,2,2,2),(3,3,3,1)\}\\
 (5,3,1,1)&\to\{(5,2,2,1),(4,4,1,1)\} &\quad (4,3,3)&\to\{(4,3,1^3),(4,2,2,2)\}
\end{array}
\]

\subsection{Two exact reductions}

\begin{proposition}[$R_1$, first row]\label{prop:R1}
If $\mu_1=\lambda_1$ then
$K_{\lambda\mu}=K_{(\lambda_2,\lambda_3,\dots),\,(\mu_2,\mu_3,\dots)}$.
\end{proposition}
\begin{proof}
In $T\in\SSYT(\lambda,\mu)$ every $1$ lies in row $1$ (a $1$ in row $i\ge2$ would need a
smaller entry above it) and, rows being weakly increasing, the $1$s occupy the leftmost
cells of row $1$.  As $\mu_1=\lambda_1$ they fill row $1$ exactly.  Deleting row $1$ is
injective into $\SSYT((\lambda_2,\dots),(\mu_2,\dots))$; conversely, prepending a row of
$\lambda_1$ ones to any such tableau is column-strict because all its entries are $\ge2$.
\end{proof}

\begin{proposition}[$R_2$, first column]\label{prop:R2}
If $\ell(\mu)=\ell(\lambda)=m$ then $K_{\lambda\mu}=K_{\hat\lambda\hat\mu}$ with
$\hat\lambda=(\lambda_1-1,\dots,\lambda_m-1)$, $\hat\mu=(\mu_1-1,\dots,\mu_m-1)$
(zero parts dropped).
\end{proposition}
\begin{proof}
Entries of $T$ lie in $\{1,\dots,\ell(\mu)\}=\{1,\dots,m\}$, and column $1$ has $m$ cells
with strictly increasing entries, so column $1$ is $(1,2,\dots,m)^{\mathsf T}$.  Deleting
it gives a semistandard tableau of shape $\hat\lambda$ and content $\hat\mu$.
Conversely, given such an $S$, prepend the column $(1,\dots,m)^{\mathsf T}$: columns stay
strict, and rows stay weakly increasing because in any column-strict tableau with
positive entries $S(i,1)\ge i$.  The two maps are mutually inverse.
\end{proof}

Note $\lambda\dom\mu$ forces $\ell(\mu)\ge\ell(\lambda)$ and $\mu_1\le\lambda_1$, so
$R_1$ and $R_2$ apply at the two opposite extremes.  Both verified on all $204$ pairs
with $N\le9$ to which they apply: $0$ failures.

\begin{corollary}\label{cor:families}
Each of the following has $K_{\lambda\mu}=1$, hence $K_{\lambda\mu}(t)$ a monomial.
\begin{enumerate}
\item[(a)] $\lambda=(N)$, any $\mu$ \ (Macdonald III \S6, Example 1: $K_{(N)\mu}(t)=t^{n(\mu)}$).
\item[(b)] $\lambda=(N-j,j)$ and $\mu=(N-k,k)$ with $j\le k$: apply $R_2$ $j$ times to
reach a one-row shape.  So the two-row/two-row slice is a singleton fibre throughout.
\item[(c)] \emph{Adjacent transfer}: $\mu=\lambda-r\,\varepsilon_k+r\,\varepsilon_{k+1}$
for some $k\ge1$, $r\ge1$, whenever this is a partition.
\end{enumerate}
\end{corollary}
\begin{proof}[Proof of (c)]
By $R_1$ applied $k-1$ times we may take $k=1$, so
$\mu=(\lambda_1-r,\lambda_2+r,\lambda_3,\dots)$.  In $T\in\SSYT(\lambda,\mu)$ the
$\lambda_1-r$ ones fill the left of row $1$.  Entries in rows $\ge3$ are $\ge3$, so all
$\mu_2=\lambda_2+r$ twos lie in rows $1$ and $2$; row $2$ holds at most $\lambda_2$ of
them and row $1$ has exactly $r$ free cells, so row $1=1^{\lambda_1-r}2^{r}$ and row
$2=2^{\lambda_2}$, forced.  (This is column-strict because $\mu_1\ge\mu_2$ gives
$\lambda_1-r\ge\lambda_2+r\ge\lambda_2$.)  The remaining letters are
$3^{\lambda_3},4^{\lambda_4},\dots$ on the shape $(\lambda_3,\lambda_4,\dots)$, where the
only semistandard filling is the one with row $i$ constant --- i.e.\ $K_{\nu\nu}=1$.
\end{proof}

Verified: all $250$ adjacent-transfer pairs with $N\le11$ have $K=1$, and all $250$ are
in fact killed by $R_1$ and $R_2$ alone.

\subsection{The irreducible core}

Call $(\lambda,\mu)$ \emph{irreducible} if $\ell(\lambda)\ge2$, $\mu_1<\lambda_1$ and
$\ell(\mu)>\ell(\lambda)$, i.e.\ neither $R_1$ nor $R_2$ applies and $\lambda$ is not a
single row.

\begin{theorem}[measured]\label{thm:core}
Of the $1044$ irreducible dominance pairs with $N\le10$, exactly $\mathbf{10}$ have
$K_{\lambda\mu}=1$:
\[
\begin{array}{lll}
 ((3,3),(2,2,2)) & ((4,4),(3,3,2)) & ((5,5),(4,4,2))\\
 ((5,5),(4,3,3)) & ((4,2,2),(3,3,1,1)) & ((5,2,2),(4,3,1,1))\\
 ((6,2,2),(5,3,1,1)) & ((6,2,2),(4,4,1,1)) & ((5,3,2),(4,4,1,1))\\
 ((4,2,2,2),(3,3,2,1,1)) & &
\end{array}
\]
\end{theorem}

A third exact symmetry, rectangular complementation --- with $r=\max(\ell(\lambda),\ell(\mu))$,
$c=\lambda_1$, $K_{\lambda\mu}=K_{\lambda^*\mu^*}$ where
$\lambda^*=(c-\lambda_r,\dots,c-\lambda_1)$, $\mu^*=(c-\mu_r,\dots,c-\mu_1)$ (the
$\mathrm{GL}_r$ duality $V_\lambda^*\otimes\det^{\,c}$) --- disposes of four of the ten,
sending them to one-row shapes:
\[
 ((3,3),(2,2,2))\mapsto((3),(1^3)),\quad ((4,4),(3,3,2))\mapsto((4),(2,1,1)),
\]
\[
 ((5,5),(4,4,2))\mapsto((5),(3,1,1)),\quad ((5,5),(4,3,3))\mapsto((5),(2,2,1)).
\]
The remaining six are not reduced by it; three of them --- $((4,2,2),(3,3,1,1))$,
$((5,3,2),(4,4,1,1))$, $((4,2,2,2),(3,3,2,1,1))$ --- are fixed points of the
complementation.  All six have $\lambda-\mu$ a sum of \emph{two disjoint} adjacent
transfers, e.g.\ $\lambda-\mu=(1,-1,1,-1)$ for the first three; but almost all pairs of
that form have $K\ge2$, so this is a description of the six, not a criterion.

\subsection{Two things the boundary is not a function of}\label{sec:obstruction}

\begin{theorem}\label{thm:obstruction}
\begin{enumerate}
\item[(a)] No criterion depending only on the difference vector $\lambda-\mu$ can decide
whether $K_{\lambda\mu}=1$.  Witness:
$\lambda-\mu=(1,1,-2)$ for both $((3,3),(2,2,2))$, where $K=1$ and $K_{\lambda\mu}(t)=t^3$,
and $((4,3),(3,2,2))$, where $K=2$ and $K_{\lambda\mu}(t)=t^2+t^3$.
\item[(b)] No criterion depending only on the dominance-slack vector
$s_k=\sum_{i\le k}(\lambda_i-\mu_i)$ can decide it.  Witness: $s=(2,1,0)$ for both
$((4,1,1),(2,2,2))$, where $K=1$ and $K_{\lambda\mu}(t)=t^3$, and $((5,1),(3,2,1))$,
where $K=2$ and $K_{\lambda\mu}(t)=t^2+t^3$.
\end{enumerate}
\end{theorem}
\begin{proof}
The four Kostka numbers and four polynomials are computed by both engines and listed
above; the two difference vectors (resp.\ slack vectors) are equal by inspection.
\end{proof}

These two are worth recording because they are exactly the two shapes of answer I would
have reached for first.  (a) kills ``$\mu$ is obtained from $\lambda$ by the following
box moves''; (b) kills ``$\lambda$ and $\mu$ touch in the following positions''.  Any
correct criterion must read $\lambda$ and $\mu$ jointly, not through either of these
projections.

\section{Gaps}\label{sec:gaps}

\begin{gap}
A closed-form characterisation of $\{(\lambda,\mu):K_{\lambda\mu}=1\}$ is \textbf{not}
obtained.  What is obtained is: the filter structure (Prop.~\ref{prop:filter}), the two
exact reductions, the proved families of Cor.~\ref{cor:families}, the measured statement
that the reductions leave only $10$ of $1044$ irreducible pairs for $N\le10$, and the
proved fact (Thm.~\ref{thm:obstruction}) that no criterion in $\lambda-\mu$ or in the
slack vector can exist.  The precise shape of the obstruction is the six-element
post-complementation core: it is nonempty, so ``$K_{\lambda\mu}=1$ iff $(\lambda,\mu)$
reduces to a one-row shape under $R_1,R_2,R_3$'' is \emph{false}.
\end{gap}

\begin{gap}
Theorem~\ref{thm:minimal} ($\le2$ minimal elements) and Theorem~\ref{thm:core}
($10$ of $1044$) are measurements for $N\le10$, with no proof and no conjecture attached.
In particular I do \emph{not} claim the bound $2$ persists; the honest reading is that
the exceptional set is sparse in this range, and sparse is not empty.
\end{gap}

\begin{gap}
Prop.~\ref{prop:filter} is cited, not proved: the $\mathfrak{sl}_2$-string decomposition
of $B(\lambda)$ is standard and I have sketched how it gives monotonicity, but I have not
written the argument out.  It is used only for the \emph{structure} of the boundary
(Thm.~\ref{thm:minimal}), never for Theorem~\ref{thm:main}.
\end{gap}

\paragraph{What the main theorem does and does not rest on.}
Theorem~\ref{thm:main} rests on Lemmas~\ref{lem:at1} and \ref{lem:monic} only, both from
Macdonald III \S6, both deep-read, and Lemma~\ref{lem:monic}'s proof is reproduced in
full.  It does \emph{not} rest on the Lascoux--Schützenberger theorem III (6.5)(i), which
enters only in Corollary~\ref{cor:charge} where the word ``charge'' is used, and not in
the computations, which agree across two mechanism-disjoint engines and with Macdonald's
printed tables.

\paragraph{Numerals appearing in the answer.}  Per the brief's \S4(f): the only constants
in the statements are $1$ (the Kostka number, the leading coefficient, and the number of
maximal-charge tableaux --- these are three different occurrences of $1$ and
Lemma~\ref{lem:monic} is exactly the bridge between the second and the third), and the
exponent $n(\mu)-n(\lambda)$, which at $\lambda=\mu$ is $0$ and must not be confused with
the charge of the superstandard tableau: they agree, but for the reason that $K_{\lambda\lambda}(t)=1$,
not by coincidence.  The $2$ of Theorem~\ref{thm:minimal} is a measured cardinality of a
set of minimal elements; it is not $\ell$ of anything, not a Bell number, and the first
place a competing reading would diverge is $N=11$, which is not computed.

\paragraph{Code.}  \texttt{proofs/code-1009-q380/}: \texttt{kf.py} (Engine A),
\texttt{kostant.py} (Engine B), \texttt{census.py}, \texttt{controls.py},
\texttt{reduce.py}.

\end{document}
